\clearpage\SourceLeaf{134}%
\pdfbookmark[1]{On a Problem of Spherical Astronomy}{sphaerisch1808}
\fancyhead[L]{\small ON A PROBLEM OF SPHERICAL ASTRONOMY}
\begin{center}
{\small ON A PROBLEM OF SPHERICAL ASTRONOMY.\par}
\medskip\hrule\medskip
{\small Monatliche Correspondenz zur Beförderung der Erd- und Himmels-Kunde,\par
edited by Baron \textsc{von Zach}. October 1808.\par}
\medskip\hrule
\end{center}\medskip
The daily motion of celestial bodies gives rise to a great variety of problems concerning the relations among hour angles, altitudes, azimuths, the positions of celestial bodies, and the elevation of the pole. \textsc{Maupertuis} and other astronomers have studied several of these problems, which vary greatly in practical usefulness and in some cases have very little. Such problems have been recommended especially to seafarers and travelers, and also to observers equipped with only a few, less refined instruments, to help with the two most essential observations, determining the time and the position. Astronomers who have the finest clocks, transit instruments, mural quadrants, full circles, and zenith sectors at their disposal need no provision of this kind. At any hour they can determine the time with the greatest precision and ease, and for them the question of which methods are best suited to establishing the geographical position of their observing site is already settled. Naturally, determining the positions of the stars and the Sun also remains the province of these astronomers alone. An observer with poorer instruments will always take those positions from them when needed and will have no wish to derive them from his own less exact observations. So, for example, the problem of determining from three observed stellar altitudes
\SourceLeaf{135}%
the elevation of the pole, the declination of the star, and the time of its culmination all at once has no practical value at all, except in the one hardly conceivable case in which there would be an opportunity to observe only a single star whose position had yet to be well determined. That procedure could give passable results only if the observations were very widely spaced in time and both the clock’s rate during them and the measured altitudes themselves were very accurate. Even then, one can always derive the elevation of the pole and the time \textit{much} more precisely from two of those observations by treating the star’s declination as given.

One of the most useful problems of all for seafarers and traveling observers is to determine the clock offset and the elevation of the pole from two observed altitudes of two stars and the corresponding clock readings. The stars’ right ascensions and declinations are treated as given, and the clock either keeps sidereal time or is assumed to have a known rate during the observations. The stars and the altitudes are entirely at one’s choice, subject only to the condition that the vertical circles in which the altitudes are measured form an angle at the zenith that is neither too acute nor too close to $180^0$. In this way the two altitudes can easily be measured within a few minutes. This is undoubtedly of the greatest importance both for the seafarer, whose position at sea is always changing rapidly, and for the observer traveling on land, whose time is limited, who may have favorable weather for only a few minutes, or who cannot rely on the rate of his clock for long. The direct solution, which readily suggests itself, rests on calculating three spherical triangles and is admittedly somewhat lengthy. The indirect method can be made very convenient and flexible. I reserve the subject for another occasion.

When observing stellar altitudes with a reflecting instrument and an artificial horizon, one finds it somewhat difficult at first to bring both images into the field of view. One soon acquires the knack, especially by keeping to the brighter stars. A very rough preliminary calculation of the altitude helps, if the elevation of the pole and the clock offset are already known approximately. If one can observe from a stand,
\SourceLeaf{136}%
it is much more convenient, and one can see the contact of the images much more precisely and closer to the center of the field of view than when observing with the instrument held in the hand. With diligent practice one can nevertheless get quite far at this too.

One inconvenience of the method just described is that the fine graduations of the sextant cannot be read as well by artificial light at night as they can by day. A more important point, which has come to attention only recently, is that many reflecting sextants, even among those made by the finest makers, have graduations that fall somewhat short of the perfection we have hitherto been accustomed to assume as a matter of course. I myself own a 10-inch reflecting sextant by \textsc{Troughton}, No.\,420. It is excellent in other respects, yet it quite decidedly reads angles from $100$ to $120^0$ too small by $50$ to $60$ or $70''$. All the other tests and adjustments have been carried out on it with the greatest care, so I can attribute this only to an error in the graduations. Such errors are of no importance when determining latitude and time at sea and of little importance on land at sites whose positions were previously quite unknown. In cases where greater precision is desired, however, one cannot take the perfection of the graduations on trust. One must first try to ascertain their errors as accurately as possible, or give up making fully reliable determinations by methods that require precisely measured altitudes.

For these reasons, many observers should welcome a method by which they can determine the elevation of the pole and the clock offset very precisely in a short time, however poorly graduated the instrument may be, or even if it has no graduations at all. In this way they can also determine any graduation errors in the instrument itself. The accuracy of the results thus depends solely on how carefully one observes the contact of the images, since the positions of the stars used may be regarded as entirely free from error in the present state of our star catalogs. The method consists in waiting for the times when three stars of one’s choice reach the same altitude in vertical circles that do not form too acute an angle at the zenith. The altitude itself can be chosen at will and need not
\SourceLeaf{137}%
be known. The following will be the most convenient way to derive the elevation of the pole and the clock offset from these observations. I use the following symbols.
\begin{description}
\item[$\alpha,\alpha',\alpha''$] the right ascensions of the three stars,
\item[$\delta,\delta',\delta''$] their declinations, with southern declinations taken as negative,
\item[$\theta,\theta',\theta''$] the three clock readings when these stars reach altitude $h$,
\item[$k$] the amount by which the clock is ahead of sidereal time, which I assume to be the same for all three observations. If the clock did not keep sidereal time, one could take $k$ to be its advance at any chosen instant. In that case $\theta,\theta',\theta''$ would mean the times the clock would have shown had it kept exact sidereal time between that instant and the observations. One might, for instance, take $k$ to be its advance at the first observation and increase or decrease the other two times by the appropriate amounts if the clock were running slower or faster than sidereal time.
\item[$\varphi$] the elevation of the pole at the site.
\end{description}
The three hour angles are evidently $\theta-k-\alpha$, $\theta'-k-\alpha'$, $\theta''-k-\alpha''$, converted from time to arc. Thus, if one sets
\[
\begin{aligned}
\theta-\alpha&=t\\
\theta'-\alpha'&=t'\\
\theta''-\alpha''&=t''
\end{aligned}
\]
one obtains the following three equations.
\[
\begin{aligned}
\mathrm{I.}\qquad&\sin h=\sin\varphi\sin\delta+\cos\varphi\cos\delta\cos(t-k)\\
\mathrm{II.}\qquad&\sin h=\sin\varphi\sin\delta'+\cos\varphi\cos\delta'\cos(t'-k)\\
\mathrm{III.}\qquad&\sin h=\sin\varphi\sin\delta''+\cos\varphi\cos\delta''\cos(t''-k)
\end{aligned}
\]

Subtracting I from II and making a simple transformation gives
\[
\begin{aligned}
&2\sin\varphi\sin\tfrac12(\delta'-\delta)\cos\tfrac12(\delta'+\delta)\\
&\quad=2\cos\varphi\cos[\tfrac12(t+t')-k]\cos\tfrac12(t'-t)\sin\tfrac12(\delta'-\delta)\sin\tfrac12(\delta'+\delta)\\
&\qquad+2\cos\varphi\sin[\tfrac12(t+t')-k]\sin\tfrac12(t'-t)\cos\tfrac12(\delta'-\delta)\cos\tfrac12(\delta'+\delta)
\end{aligned}
\]
or
\[
\begin{aligned}
\operatorname{tang}\varphi&=\cos[\tfrac12(t+t')-k]\cos\tfrac12(t'-t)\operatorname{tang}\tfrac12(\delta'+\delta)\\
&\quad+\sin[\tfrac12(t+t')-k]\sin\tfrac12(t'-t)\operatorname{cotg}\tfrac12(\delta'-\delta)
\end{aligned}
\]
\SourceLeaf{138}%
If one determines $A'$ and $B'$ so that
\[
\begin{aligned}
A'\sin B'&=\sin\tfrac12(t'-t)\operatorname{cotg}\tfrac12(\delta'-\delta)\\
A'\cos B'&=\cos\tfrac12(t'-t)\operatorname{tang}\tfrac12(\delta'+\delta)
\end{aligned}
\]
and also sets
\[
\tfrac12(t+t')-B'=C'
\]
the preceding equation takes the form
\[
\mathrm{IV.}\qquad\operatorname{tang}\varphi=A'\cos(C'-k)
\]

In exactly the same way, or simply by exchanging the quantities relating to the second observation with those relating to the third, one sees that if $A''$ and $B''$ are determined so that
\[
\begin{aligned}
A''\sin B''&=\sin\tfrac12(t''-t)\operatorname{cotg}\tfrac12(\delta''-\delta)\\
A''\cos B''&=\cos\tfrac12(t''-t)\operatorname{tang}\tfrac12(\delta''+\delta)
\end{aligned}
\]
and if one also sets
\[
\tfrac12(t+t'')-B''=C''
\]
the following equation results
\[
\mathrm{V.}\qquad\operatorname{tang}\varphi=A''\cos(C''-k)
\]

It is now easy to derive $k$ and $\varphi$ by combining IV and V. The equations give
\[
A'\cos(C'-k)=A''\cos(C''-k)
\]
hence
\[
\begin{aligned}
&(A''-A')\cos[\tfrac12(C'+C'')-k]\cos\tfrac12(C''-C')\\
&=(A'+A'')\sin[\tfrac12(C'+C'')-k]\sin\tfrac12(C''-C')
\end{aligned}
\]
Thus, if one sets
\[
\frac{A'}{A''}=\operatorname{tang}\zeta
\]
which gives
\[
\frac{A''-A'}{A''+A'}=\operatorname{tang}(45^0-\zeta)
\]
\SourceLeaf{139}%
and determines $\psi$ from the equation
\[
\operatorname{tang}(45^0-\zeta)\operatorname{cotang}\tfrac12(C''-C')=\operatorname{tang}\psi
\]
one obtains
\[
k=\tfrac12(C'+C'')-\psi
\]
and then $\varphi$ from either equation IV or equation V. From the resulting elevation of the pole and the hour angles $t-k$, $t'-k$, $t''-k$, one can then, if desired, derive the altitude $h$ by formulas I, II, III or by other known methods. Its value must be the same from all three observations.

The following remarks should be added to this solution.

1) From the two equations $A\sin B=M$, $A\cos B=N$, one determines $A$ and $B$ using the formulas
\[
\begin{gathered}
\operatorname{tang}B=\frac{M}{N}\\
A=\frac{M}{\sin B}\quad\text{or}\quad A=\frac{N}{\cos B}
\end{gathered}
\]
The first expression for $A$ is preferred when $M$ is greater than $N$, and the other in the opposite case. When determining the angle $B$ from its tangent, one is free to choose between the first and third quadrants if the tangent is positive, or between the second and fourth if it is negative. One chooses the quadrant in which the sine has the sign of $M$ and the cosine the sign of $N$, which makes $A$ positive.

2) The similar ambiguity in determining $\psi$ from the tangent must be resolved in such a way that $\operatorname{tang}\varphi$ is positive. One therefore takes $\psi$ between $-90^0$ and $+90^0$, assuming, of course, that the observations were made in the northern hemisphere. In the southern hemisphere it would be the other way around.

For every method proposed for use by practical astronomers, it is an indispensable duty to assess the effect of unavoidable observational errors on the results. That is how one can tell whether the method can be used reliably at all, and under what circumstances. Neglect of this duty accounts for the many half-baked ideas whose uselessness practical astronomers complain about. A simple investigation, omitted here for brevity,
\SourceLeaf{140}%
shows the following. Let $\lambda,\lambda',\lambda''$ be the azimuths corresponding to the three observations, reckoned from the south point in the direction of the daily motion. An error of $\Delta$ seconds of time in the first contact then entails an error of
\[
\frac{\Delta\sin\lambda\sin\tfrac12(\lambda''+\lambda')}{2\sin\tfrac12(\lambda'-\lambda)\sin\tfrac12(\lambda''-\lambda)}\quad\text{seconds of time}
\]
in the determination of the time and of
\[
\frac{15\Delta\sin\lambda\cos\varphi\cos\tfrac12(\lambda''+\lambda')}{2\sin\tfrac12(\lambda'-\lambda)\sin\tfrac12(\lambda''-\lambda)}\quad\text{seconds of arc}
\]
in the elevation of the pole. Indeed, if the contact was observed too early, $k$ must be decreased and $\varphi$ increased by these amounts. The expressions for the effect of errors in the second and third observations are quite similar and are obtained simply by exchanging $\lambda$ with $\lambda'$ or $\lambda''$. Since the altitude of the first star decreases in $\Delta$ seconds of time by $15\Delta\sin\lambda\cos\varphi$ seconds of arc, the effect of the observational errors can be expressed even more simply. Suppose that at the instant when the first star’s altitude was found to agree with the setting of the instrument, it was in fact still or already $D$ seconds of arc higher. The value obtained for $k$, in seconds of arc, is then too large by
\[
\frac{D\sin\tfrac12(\lambda''+\lambda')}{2\cos\varphi\sin\tfrac12(\lambda'-\lambda)\sin\tfrac12(\lambda''-\lambda)}
\]
and that for $\varphi$ too small by
\[
\frac{D\cos\tfrac12(\lambda''+\lambda)}{2\sin\tfrac12(\lambda'-\lambda)\sin\tfrac12(\lambda''-\lambda)}
\]

It follows that one need only make sure that none of the sines of $\tfrac12(\lambda'-\lambda)$, $\tfrac12(\lambda''-\lambda)$, $\tfrac12(\lambda''-\lambda')$ becomes too small. This is done by choosing only stars that reach altitude $h$ at fairly different azimuths. Second, it is clear that stars whose altitudes change slowly are just as useful as those that rise or fall rapidly. With the former there is no need to pinpoint the instant when they reach the required altitude. All that matters is that at the instant taken for it, they are in fact at no appreciable distance from that altitude. One can therefore choose stars near culmination or the Pole Star without hesitation. Such stars are especially suitable,
\SourceLeaf{141}%
because there is no hurry in meeting the requirement just mentioned. At least one of the three stars will in any case always change its altitude more quickly if the condition of differing azimuths is satisfied.

It is good to prepare before making the observation. At any hour one will easily find some recognizable stars that reach the same altitude in quick succession. A globe or a stereographic projection of the celestial sphere is all one needs to do this easily. The observation will go all the better if one knows beforehand, to within a couple of minutes, the intervals at which the three observations follow one another. One would have to have observed poorly if this failed to give the clock offset to within a couple of seconds at least. By combining several stars and repeating the observation on several evenings, when the same altitudes can be measured again, one will be able to achieve a very high degree of accuracy with this method. There are also several ways to shorten the calculation, particularly when the same stars have often been observed at the same or nearly the same altitude. I do not go into these here.

To illustrate this method further, I give the full calculation for an actual application. On 27 August of this year I used a mercury horizon and set my sextant to twice the altitude, $105^0\ 18'\ 55''$. Holding the sextant in my hand, I observed the stars $\alpha$ Andromedae, $\alpha$ Ursae Minoris, and $\alpha$ Lyrae, using the \textsc{Shelton} pendulum clock, which ran at the exact sidereal rate. The readings were as follows.
\begin{center}
\begin{tabular}{lrrr}
$\alpha$ Andromedae $\ .\ .\ .\ .$&$21^{\mathrm u}$&$33^{\mathrm m}$&$26^{\mathrm s}$\\
$\alpha$ Ursae Minoris $\ .\ .\ .\ .$&$21$&$47$&$30$\\
$\alpha$ Lyrae $\ .\ .\ .\ .\ .\ .$&$22$&$5$&$21$
\end{tabular}
\end{center}
The apparent positions of the stars on those days were as follows.
\begin{center}
\small
\begin{tabular}{lcc}
&Right ascension&Northern declination\\
$\alpha$ Andromedae&$23^{\mathrm u}\ 58^{\mathrm m}\ 33^{\mathrm s}33$&$28^0\ 2'\ 14''8$\\
$\alpha$ Ursae Minoris&$0\quad55\quad4.7$&$88\quad17\quad5.7$\\
$\alpha$ Lyrae $\ .\ .$&$18\quad30\quad28.96$&$38\quad37\quad6.6$
\end{tabular}
\end{center}
From these we obtain the following.
\SourceLeaf{142}%
\begin{center}\small
\begin{tabular}{rcrcl}
&In time&&&In arc\\
$t=$&$21^{\mathrm u}\ 34^{\mathrm m}\ 52^{\mathrm s}67$&$\ .\ .\ .$&&$323^0\ 43'\ 10''05$\\
$t'=$&$20\quad52\quad25.30$&$\ .\ .\ .$&&$313\quad6\quad19.50$\\
$t''=$&$3\quad34\quad52.04$&$\ .\ .\ .$&&$53\quad43\quad0.60$
\end{tabular}
\par\medskip
\begin{minipage}[t]{.49\textwidth}\centering
\begin{tabular}[t]{@{}lr@{}}
$\tfrac12(t'-t)=$&$-5^0\ 18'\ 25''27$\\
$\tfrac12(t'+t)=$&$318\quad24\quad44.77$\\[6pt]
$\tfrac12(\delta'-\delta)=$&$30\quad7\quad25.45$\\
$\tfrac12(\delta'+\delta)=$&$58\quad9\quad40.25$\\[6pt]
$\log\sin\tfrac12(t'-t)\;.\;.$&$8.9661069\,\mathrm n$\\
$\log\cot\tfrac12(\delta'-\delta)\;.\;.$&$0.2363974$\\[6pt]
$\log\cos\tfrac12(t'-t)\;.\;.$&$9.9981343$\\
$\log\operatorname{tang}\tfrac12(\delta'+\delta)\;.\;.$&$0.2069331$
\end{tabular}
\end{minipage}\hfill
\begin{minipage}[t]{.49\textwidth}\centering
\begin{tabular}[t]{@{}lr@{}}
$\tfrac12(t''-t)=$&$-135^0\ 0'\ 4''72$\\
$\tfrac12(t''+t)=$&$188\quad43\quad5.32$\\[6pt]
$\tfrac12(\delta''-\delta)=$&$5\quad17\quad25.90$\\
$\tfrac12(\delta''+\delta)=$&$33\quad19\quad40.70$\\[6pt]
$\log\sin\tfrac12(t''-t)\;.\;.$&$9.8494751\,\mathrm n$\\
$\log\cot\tfrac12(\delta''-\delta)\;.\;.$&$1.0333869$\\[6pt]
$\log\cos\tfrac12(t''-t)\;.\;.$&$9.8494949\,\mathrm n$\\
$\log\operatorname{tang}\tfrac12(\delta''+\delta)\;.\;.$&$9.8179461$
\end{tabular}
\end{minipage}
\end{center}
We therefore have
\begin{center}\small
\begin{minipage}[t]{.49\textwidth}\centering
\begin{tabular}[t]{@{}lr@{}}
$\log A'\cos B'\;.\;.\;.$&$0.2050674$\\
$\log A'\sin B'\;.\;.\;.$&$9.2025043\,\mathrm n$
\end{tabular}
\end{minipage}\hfill
\begin{minipage}[t]{.49\textwidth}\centering
\begin{tabular}[t]{@{}lr@{}}
$\log A''\cos B''\;.\;.\;.$&$9.6674410\,\mathrm n$\\
$\log A''\sin B'\;.\;.\;.$&$0.8828620\,\mathrm n$
\end{tabular}
\end{minipage}
\end{center}
From which we obtain
\begin{center}\small
\begin{minipage}[t]{.49\textwidth}\centering
\begin{tabular}[t]{@{}lr@{}}
$B'=$&$354^0\ 19'\ 22''04$\\
$B''=$&$266\quad30\quad55.07$\\[6pt]
$C'=$&$-35\quad54\quad37.27$\\
$C''=$&$-77\quad47\quad49.75$\\[6pt]
$\tfrac12(C''-C')=$&$-20^0\ 56'\ 36''24$\\
$\tfrac12(C''+C')=$&$-56\quad51\quad13.51$
\end{tabular}
\end{minipage}\hfill
\begin{minipage}[t]{.49\textwidth}\centering
\begin{tabular}[t]{@{}lr@{}}
$\log A'\;.\;.\;.$&$0.2072029$\\
$\log A''\;.\;.\;.$&$0.8836657$\\[6pt]
$\log\operatorname{tang}\zeta\;.\;.$&$9.3235372$\\
$\zeta=$&$11^0\ 53'\ 41''28$\\[6pt]
$45^0-\zeta=$&$33\quad6\quad18.72$
\end{tabular}
\end{minipage}
\par\medskip
\begin{tabular}{lr}
$\log\operatorname{tang}(45^0-\zeta)\;.\;.$&$9.8142617$\\
$\log\cot\tfrac12(C''-C')\;.\;.$&$0.4171063\,\mathrm n$\\\hline
$\log\operatorname{tang}\psi\;.\;.\;.\;.$&$0.2313680\,\mathrm n$\\
\multicolumn{2}{l}{$\psi=-59^0\ 35'\ 14''71$}\\
\multicolumn{2}{l}{$k=+2\quad44\quad1.20$}
\end{tabular}
\end{center}
The clock was therefore ahead of sidereal time by $10^{\mathrm m}\,56^{\mathrm s}08$.
\SourceLeaf{143}%
This agrees within a couple of tenths with the value derived from stellar transits observed with the mural quadrant.

To determine the elevation of the pole, we have
\begin{center}\small
\begin{tabular}{@{}lr@{}p{.08\linewidth}@{}lr@{}}
\multicolumn{2}{@{}l@{}}{$C'-k=-38^0\ 38'\ 38''47$}&\hspace{.07\linewidth}&&\\
\multicolumn{2}{@{}l@{}}{$C''-k=-80\quad31\quad50.95$}&&&\\
$\log A'\;.\;.\;.\;.$&$0.2072029$&&$\log A''\;.\;.\;.\;.$&$0.8836657$\\
$\log\cos(C'-k)\;.$&$9.8926738$&&$\log\cos(C''-k)\;.$&$9.2162109$\\
\cline{1-2}\cline{4-5}
$\log\operatorname{tang}\varphi\;.\;.\;.$&$0.0998767$&&&$0.0998766$\\
\multicolumn{2}{@{}l@{}}{$\varphi=51^0\ 31'\ 51''51$}&&&
\end{tabular}
\end{center}

Using these results to calculate the true altitudes of the stars gives the same value from all three to within a hundredth of a second.
\[
h=52^0\ 37'\ 21''3
\]
Refraction, corrected for the barometer and thermometer readings, is $42''7$, so the apparent altitude is $52^0\ 38'\ 4''$. The sextant’s collimation error is $-3'\ 30''$, so the measured angle is $105^0\ 15'\ 25''$. Since the true angle is $105^0\ 16'\ 8''$, the measured angle is too small by $43''$.

On 25 August the sextant was set to an angle $5''$ larger, namely $105^0\ 19'\ 0''$. The same three stars were observed at the clock readings below.
\begin{center}
\begin{tabular}{lrrr}
$\alpha$ Andromedae $\ .\ .\ .\ .$&$21^{\mathrm u}$&$33^{\mathrm m}$&$29^{\mathrm s}$\\
$\alpha$ Ursae Minoris $\ .\ .\ .\ .$&$21$&$47$&$38$\\
$\alpha$ Lyrae $\ .\ .\ .\ .\ .\ .$&$22$&$5$&$22$
\end{tabular}
\end{center}
This gives the following results.
\begin{center}
\begin{tabular}{lr}
Clock advance over sidereal time&$10^{\mathrm m}\,57^{\mathrm s}9$\\
$\varphi=$&$51^0\ 31'\ 56''7$\\
$h=$&$52\quad37\quad29.15$
\end{tabular}
\end{center}
Twice the apparent altitude is $105^0\ 16'\ 25''5$. The sextant gives $105^0\ 15'\ 30''$, so its error is $-54''5$.

On average, then, the sextant error at the angle $105^0$ is $-48''7$, and the elevation of the pole is $51^0\ 31'\ 54''1$. Some observations by Professor \textsc{Harding}
\SourceLeaf{144}%
of the same stars, but at a somewhat greater altitude, and an earlier observation of my own with $\alpha$ Andromedae, $\alpha$ Aquilae, and $\alpha$ Lyrae, give the same result for the elevation of the pole within a couple of seconds.

To see the effect of the observational errors more clearly, I used the numbers of the example worked out here to calculate the azimuths corresponding to the three stars. I obtained the values below.
\begin{center}
\begin{tabular}{lcl}
$\alpha$ Andromedae & $.\ .\ .\ .$ & $293^{0}\ 45'\ 15''$\\
$\alpha$ Ursae Minoris & $.\ .\ .\ .$ & $182\quad9\quad9$\\
$\alpha$ Lyrae & $.\ .\ .\ .\ .\ .$ & $90\quad17\quad52$
\end{tabular}
\end{center}
From this I find, using the formulas above, that if the three contacts were observed $\Delta,\Delta',\Delta''$ seconds of time too early, the computed advance of the clock is too small by
\[
+0.391\Delta+0.0066\Delta'+0.603\Delta''
\]
seconds of time, and the computed elevation of the pole is too small by
\[
+3.808\Delta-0.2884\Delta'-3.519\Delta''
\]
seconds of arc.

Thus if the observation of $\alpha$ Ursae Minoris is off by $20''$, and that of each of the other two stars by $1''$, the error in the time determination can be at most $1^{\mathrm{s}}1$, and that in the elevation of the pole at most $13''$, provided the positions of the stars themselves are exactly reliable. The other formulas given above serve to determine the possible errors in the final results from errors in the altitudes instead of errors in the times. Let us suppose that at the three instants of observation the altitudes of the stars were $D,D',D''$ seconds of arc greater than required by the setting of the sextant. This gives an error of
\[
+0.0430D+0.0177D'-0.0607D''
\]
seconds of time in the advance of the clock, and of
\[
+0.446D-0{,}823D'+0.377D''
\]
seconds of arc in the elevation of the pole. Both values come out too large by these respective amounts. The latter formulas also serve to estimate the effect of any remaining uncertainty in the positions of the stars.
\SourceLeaf{145}%
One need only take $D,D',D''$ to be the errors in these positions in the directions of the three vertical circles, in seconds of arc along a great circle, to apply the formulas at once. For $D,D',D''$ are the differences between the altitudes corresponding to the sextant setting and the altitudes of the points on the celestial sphere with the right ascensions and declinations used in the calculation. It obviously makes no difference whether they arise from observational errors when the stars coincide exactly with these points, or from small differences between the stars and those points when the observations are perfectly accurate. Finally, I note that the sum of the coefficients is always $=0$ in each of the two formulas for the error in the elevation of the pole, as well as in the second formula for the error in the determination of the time. In the first formula for the error in the determination of the time it is $=1$. With a little thought one will easily see why.

\par\bigskip\centerline{\rule{.27\linewidth}{.4pt}}
